\subsection{Study}
Studying at our faculty might seem challenging but if you follow the recommended
course of study, it might even become enjoyable.
All information about the obligatory courses, the recommended courses of study
and other duties can be found on the faculty's website.
A number of students unfortunately does not even make it through the first two
semesters of study but if you manage to overcome these difficult beginnings, you
will have a very good chance to successfully complete your studies.
Your older colleagues and teachers will give you a lot of useful advice but you
will not follow it and the night before an exam you will tell yourself: \uv{Why
did not I study more during the semester?}


\subsubsection{Study Programmes}

Our faculty offers bachelor’s, master’s and doctoral degree programmes, given in
either Czech or English. You usually start with the bachelor's programme and end
with completing the master's programme, however some of the students choose to
continue and study in doctoral degree programmes. This Cookbook will only
mention the bachelor's and master's degree programmes and exceptions for Erasmus
and exchange students.

\subsubsubsection{Bachelor's Programme}
If you are a regular student of the English programme, you might become a
bachelor of Computer Science.
The study usually takes three years and you have to pay for it. Studies are
concluded with a state final examination. Necessary conditions for taking the
State Final Exam include passing all obligatory courses, obtaining the required
number of credits for elective courses, reaching a total of at least 180 credits
and submitting a completed thesis (for the thesis defence). However, it is
possible to take the State Final Exam even without submitting the thesis and
complete it during the following year.
If you manage to pass, congratulations, you are a bachelor (Bc.).
If you are an Erasmus student, you might either study Computer Science or
Physics.

\subsubsubsection{Master's Programme}
If you are a regular student of the English programme, you might choose between
becoming the master of Computer Science or the master of Mathematics. The
studies usually take two years and if you manage to pass the State Final Exam
(the necessary conditions also include passing all obligatory courses, obtaining
the required number of credits for elective courses, submitting a completed
thesis and reaching a total of at least 120 credits), congratulations, you are a
master (Mgr.).


\subsubsection{Study Programmes}
There are three study programmes offered for the students of the English
programmes (this does not concern the Erasmus students) -- bachelor of Computer
Science, master of Computer Science and master of Mathematics. The Erasmus
students might also participate in the study programmes for future teachers and
Physics.


\subsubsection{Courses}
A wide range of courses is offered by the faculty, you can find them and their
descriptions in the SIS.

\subsubsubsection{ECTS Credits}
The ECTS credits are awarded at the end of the semester if you manage to pass
the course.
It is important to have enough credits because otherwise you might not be
allowed to continue to the next section of the study or you might not
successfully complete your Erasmus programme.
The number of credits you might obtain after mastering the course is given and
cannot be changed.
Some of the courses (differs according to the programmes) are obligatory (you
have to complete these before the state final examination), some of them are
elective (you have to complete the prescribed part of these courses before the
state final examination) and some of them are optional (all courses other than
the obligatory and elective that are offered at the university are considered as
Optional courses for the corresponding curriculum; it is up to you whether you
decide to take some of these courses).

You need to attain in total at least the minimum number of 45 ECTS credits to
enrol in the next year (so if you manage to attain over 90 credits in your first
year, you do not even have to do anything in the second year). However, the
normal number of credits required for the registration in the next year is 60
(corresponding to the sum of the credits in the recommended course of study).
The credits obtained for the completion of obligatory and elective courses must
prevail those obtained for the mastery of optional courses -- for the first
three years of the bachelor's programme the credits for optional courses are
counted up to the limit of 15 \% for the purpose of the Annual Evaluation of
Study.
The normal and minimum numbers of credits required for registration in the next
study section might be found in Rules of Study at the Faculty of Mathematics and
Physics and Code of Study and Examination of Charles University, both these
documents and other useful ones might be found on the faculty website.
\textbf{"The Annual Evaluation of study"} monitores the progress at the end of
each study section. If you are a bachelor student in the first year, you should
make the Annual Evaluation of study also after first study section (after first
semester), when you need to obtain at least 12 course credits.
And do not be mistaken -- you will have to complete your obligatory and elective
courses even though you otherwise might have enough credits to continue in your
study.
180 ECTS credits in total are required to complete the bachelor's programme (if
you study for more than 5 years, the number of credits rises due to the minimum
number of credits for every study section).

If you are an Erasmus or exchange student, you do not need to worry about the
aforementioned information. But attention, if you do not succeed in obtaining
the number of credits set by your home university, you might have to pay back
the Erasmus scholarship (or at least a part of it).

If you will not be sure about something, do not hesitate to contact SKAS or
Spolek Matfyzák, you can use e.g. this email address \url{sos@matfyzak.cz}.

\subsubsubsection{Prerequisites and others}
Course enrolment may be restricted by certain conditions (requisites), of which
the most common are the following: prerequisite (a prerequisite to Course X is a
course that must be successfully completed before you can enrol in Course X),
corequisite (a corequisite to Course X is a course that you have to enrol in at
the same time as Course X, or that you have already successfully completed) and
prohibited combination or incompatibility (courses X and Y are a prohibited
combination if it is impossible to enrol in Course X when you have already
completed, or you enrol in, Course Y).
In some cases, it is specified that completion of Course Y is equivalent, with
respect to the requirements of the curriculum, to completion of Course X; these
two courses are called equivalent or interchangeable.

If you are an Erasmus or exchange student, you do not have to worry about the
prerequisites and corequisites -- your Learning Agreement has been already
approved before your arrival to Prague based on your previous study results. And
sometimes it has not been easy to create a suitable Learning Agreement.

\subsubsubsection{Exams}
You have up to three attempts to pass an exam (provided there are still dates
available) but we strongly recommend preparing as well as you can for the first
attempt. If you do not succeed in passing the exam or obtaining the course
credit for a course, you are allowed to take the course again in the next
section of study, but a course can be repeated once at most. This information
however only applies to the regular students of the English programmes and not
to Erasmus students and exchange students.

There are four possible outcomes for an exam (the corresponding numerical values
and Czech equivalents are given in parentheses): Excellent (1 - Výborně), Very
good (2 - Velmi dobře), Good (3 - Dobře), Fail (4 - Nevyhověl). You pass an exam
if you obtain a grade of Excellent, Very good or Good; otherwise you fail.
The grades of Excellent and Very good have one advantage -- if you will need to
once again enrol at our university, you do not have to repeat these courses
(although that should not bother you in the best case).
How such an exam looks like? Exams are taken during the examination period at
the end of the semester and can be oral, written, or a combination of both
mentioned possibilities.
During the oral exam you usually have to choose a paper with the assignment,
then you have time to think and write down your solutions, and then talk about
it with the teacher. He might even let you to correct your mistakes. Sometimes
the exam takes even 8 hours.
If you are an Erasmus or exchange student, you will find the grade scale for the
conversion of your grades (so your courses could be successfully recognized at
your home university) at the back side of your official Transcript of Records
that you will get after completing all of your exams. It will also contain the
signature of the Vice-Dean for Students Affairs.

The examiners announce the exam dates via the SIS so it might be beneficial to
enable emails notifying you about the recently announced dates.
If none of the exam dates is convenient for you, you did not pass the exam at
the first attempt or you did not manage to pass it in time, do not be afraid to
contact the teacher and ask him to announce another exam date.
Most of the teachers are forthcoming and will allow you to try to pass your exam
e.g. at the end of July.

\subsubsubsection{Course credits and others}
Not every course has to be concluded with an exam, although most of the
obligatory ones are. To be even able to take an exam you usually need to obtain
a course credit.
A course credit (usually for classes) is awarded at the end of the semester. The
conditions for obtaining a course credit differ according to the course and
conditions defined by the teacher, for example involving the completion of a
test, programming an application, or writing a survey, and are specified by the
teacher at the beginning of the semester. The possible outcomes are Pass (in
Czech Započteno - Z) and Fail (Nezapočteno - K).
Erasmus students, pay attention to the courses that are not concluded with an
exam and agree with your home university in advance how such a course will be
recognized. Some of the partner universities award the gain of a course credit
(without a grade) with the worst grade possible.

\subsubsubsection{PE and languages}
It is great if you manage to learn at least a few Czech words or sentences that
will help you with the everyday communication during shopping, eating out or
travelling. And Czechs will be always pleased that you try to speak their
language.
Czech is not an easy language but the patient lecturers will help you to
familiarize yourself with at least the basics.
Apart from the Czech language you can also study e.g. German, Spanish or French
for free at Matfyz. Note to the Erasmus students: these courses are not awarded
with a grade, you will just obtain a course credit (and no grade).
And if these options are not sufficient for you, you can study for example
Korean, Portuguese or Russian at the Faculty of Arts at Charles University. Also
for free.

Concerning the physical education, Matfyz offers a wide range of PE classes
(from swimming and soccer to canoeing). It is said that the traditional sport of
Matfyz is volleyball -- there are regularly organized tournaments for for both
-- students and professors.

\subsubsection{Recommended course of study}
You do not need to follow any set course of study if you do not want to. To
enrol into the next section of study you just have to obtain enough ECTS credits
(in the right ratio of the credits attained for the completion of obligatory,
elective and optional courses) and complete all of the obligatory courses till
the end of your studies.
However, sometimes it is difficult to make sense of all of the prerequisites and
corequisites so the faculty has prepared the recommended course of study to help
you with that.
The recommended course of study is not binding, however it is advisable to
follow it. If you do so, you should be able to fulfil the requisites (see below)
considering the relations between the courses. It is also the best way to
achieve the timely graduation.
Every degree programme has its recommended course of study and you will find it
either in the SIS, in the Study guide or on the website. This information only
applies to the regular students of the English programme.

It is probably best to follow the recommended course of study, at least in the
first year. You might even enrol in some optional courses from the higher years
that you do not need to successfully complete (there are no sanctions for not
completing the optional courses, except for the possibility of not gaining the
scholarship).


\subsubsection{Student Affairs Department}
Your student advisors deal with all of your applications and issues -- the
joyful ones (such as successfully completing your studies) as well as the sad
ones (such as the exclusion).
%The students of each programme have one officer that takes care of all their
problems, the first year students have their own officers (so please be nice to
her).
The student advisors are also Erasmus coordinators for incoming students and
they take care of the other exchange students as well. Please bear this in your
mind when requesting information which can be easily found on the faculty
website.

Usually you will not encounter any queues at this department, with the exception
of some important deadlines. If you ever wanted to complain about the Student
Affairs Department, remember that Student Affairs Department at Matfyz is one of
the most pleasant compared to the other faculties.


\subsubsection{SIS}
The Study Information System (SIS) is most of the times a functional system that
decreases the need of the student and the Student Affairs Department
communicating in person.
The whole university uses it so if you are angry with it, just think about how
difficult it must be for the students of philosophy to deal with it and you
should immediately feel better. The link is \url{is.cuni.cz/studium/eng/}.

The SIS accompanies you during your whole studies. You have probably encountered
it when you submitted your study application.
You will also use it to enrol in courses, look up your timetable or while
ignoring the fact that some of your courses overlap (the obligatory courses
never overlap if you follow the recommended course of study, however it is not
possible to prevent the overlapping of optional courses).

You will also use it to sign up for the exams and then check your results. You
can also find there following information: about teachers, classes and their
locations and their language of instruction etc.


\subsubsection{Scholarships}
You might obtain a scholarship for outstanding academic achievement. Another
types of scholarships are different types of grants and the scholarship for RDI
(Research, Development, and Innovation) activities in accordance with special
legislation.
It is recommended to put your bank account number to the SIS at the beginning of
your study so the scholarships might be easily paid to you. The opportunity for
getting these scholarships applies only to the regular students of the English
programme.

\subsubsubsection{Scholarship for outstanding academic achievement}
The scholarship for outstanding achievement is awarded at the end of the
academic year; in the first year it is awarded also for the first study section
which means based on the exam results from the winter semester -- if you will
study hard from the start, you might get an allowance for the holidays.
The Dean decides how high this scholarship will be each year and it is awarded
according to two criteria -- the year and the average result of study (it also
contains the results of exams which you failed on your first try).
Attaining at least the normal number of credits is one of the necessary
conditions for obtaining a scholarship for excellent study achievement but
usually it is awarded to students who obtained at least 60 ECTS credits (30 in
the first study section), their actual period of current study has not exceeded
the standard period of study (3 years for the bachelor's programme) and their
weighted average grade achieved has not exceeded the set limit.
The detailed rules for awarding the scholarship can be found in Rules for
Awarding Scholarships and Bursaries at the Faculty of Mathematics and Physics.

\subsubsubsection{Scholarship for Outstanding RDI (Research, Development, and
Innovation), Artistic, or Other Creative Achievement Contributing to Greater
Knowledge
(\uv{the Scholarship for Outstanding Achievement})}
The scholarship for outstanding achievement
The Scholarship for Outstanding Achievement may be awarded, in compliance with
the conditions set out in the Rules, to students who have achieved, or
participated in the achievement of, an outstanding scientific result. An
outstanding scientific result means in particular publication of a research
paper or giving a presentation at a scientific conference.

\subsubsubsection{Grants}
Finally we will mention the special form of bursaries in cases worthy of special
consideration. It is a bursary for student projects funded from Faculty Student
Grants (\uv{the FSG}); the Dean decides on the awarding thereof twice a year,
once in autumn and once in spring.
Any student or group of students of full-time bachelor’s or master’s programme
of study at the Faculty can apply for the FSG bursary. An applicant may file no
more than one proposal for the FSG bursary in a given period (autumn or spring);
the same applies even if the student is a member of a group of applicants.

\subsubsubsection{Scholarship for RDI (Research, Development, and Innovation)
Activities in Accordance with Special Legislation and Doctoral bursary}
While studying in the master's programme and the doctoral study, you will have
more options to get the scholarships -- various scholarships for RDI activities
in accordance with special legislation or even state grants (for details write
to the Research and International Affairs Department
(\url{mensikova@dekanat.mff.cuni.cz}) or visit the officers personally, room
next to the office of your Student Advisor -- as the website about research and
grant activities is available only in Czech in this moment)\dots and of course,
the PhD students get their \uv{salary} in the form of doctoral bursary.